\documentclass[11pt]{article}
\usepackage[margin=1in]{geometry}
\usepackage{times}
\usepackage{hyperref}
\usepackage{enumitem}
\usepackage{titlesec}

\titleformat{\section}{\large\bfseries}{\thesection}{1em}{}
\titleformat{\subsection}{\normalsize\bfseries}{\thesubsection}{1em}{}

\title{\textbf{CampusSync: A Cross-College Collaboration Platform \\ for Hackathon Team Formation and Project Collaboration}\\[4pt]
\large Project Proposal for UCS503P}
\author{
Submitted to: Prof. Asif \\
Thapar Institute of Engineering and Technology \\[10pt]
Aadita Garg$^{*}$, CSED \quad Ridhima Goel$^{\dagger}$, CSED \quad Vedika Gupta$^{\ddagger}$, CSED
}
\date{\today}

\begin{document}
\maketitle
\let\thefootnote\relax\footnotetext{
$^{*}$Roll No: 1024030461, Email: agarg9\_be24@thapar.edu \\
$^{\dagger}$Roll No: 1024030466, Email: Rgoel\_be24@thapar.edu \\
$^{\ddagger}$Roll No: 1024030337, Email: vgupta9\_be24@thapar.edu}

\section{Higher-order goal \textit{...secondary framing}}
This project aims to make cross-institutional student collaboration as easy as within-college
collaboration, by reducing the friction students face when assembling a team for a hackathon
or a personal/academic project. While the broader motivation is to strengthen the student
developer ecosystem across colleges, the immediate engineering objective is to translate
skill-based teammate discovery into a measurable, deployable software solution.

\section{Time-to-value \textit{...primary heuristic}}
\begin{itemize}[leftmargin=*]
  \item We will deliver meaningful increments quickly by prototyping the core profile-and-match
  workflow and validating it with a small pilot group of students within a short iteration window.
  \item We will prioritize fast feedback over perfection by using weekly iterations (and targeting
  CI-backed automation early) to reduce release friction.
  \item Success is defined by what can be delivered and measured in the first iteration -- a
  working profile, listing, and match flow -- then improved based on user feedback.
\end{itemize}

\section{Problem Statement}
Hackathons and academic/personal software projects increasingly require diverse, multidisciplinary
skill sets that a single college's student body often cannot fully provide. A student strong in
one area (e.g., machine learning) may need a teammate skilled in another (e.g., frontend
development or UI/UX), but no such teammate may exist within their own institute. Common pain
points include:
\begin{itemize}[leftmargin=*]
  \item \textbf{Fragmentation:} teammate search currently happens through scattered personal
  contacts, college-specific WhatsApp/Discord groups, or generic social posts, with no shared,
  searchable channel across institutes.
  \item \textbf{No cross-college visibility:} a student at one college who needs a Flutter
  developer, and a student at another college who has that exact skill and is looking for a
  project, have no platform through which to find each other.
  \item \textbf{Low trust / verification:} open forums and generic social platforms cannot
  confirm that a profile genuinely belongs to a verified college student.
  \item \textbf{Poor skill discoverability:} even within informal groups, there is no structured
  way to search or filter people by tech stack, availability, or complementary (rather than
  identical) skills.
\end{itemize}
\textbf{Why this matters now (contextual relevance):} With hackathon participation and
inter-college technical events growing rapidly, even a lightweight, trustworthy matchmaking
layer can noticeably reduce the time students spend searching for teammates and improve the
quality of teams that get formed.

\subsection{Objectives}
\begin{itemize}[leftmargin=*]
  \item Develop a verified cross-college platform for student profiles and project/team listings.
  \item Enable students to discover teammates based on complementary skills and interests.
  \item Implement a ranked matching mechanism for project-specific requirements.
  \item Evaluate the system using time-to-first-match, cross-college match rate, and
  connection-to-team conversion.
  \item Deploy and validate the system through a small multi-college pilot.
\end{itemize}

\section{Proposed Solution \textit{...Software Proposition}}
\subsection{Overview}
Build a web-based ``teammate-to-team'' platform, \textbf{CampusSync}, with the following
components:
\begin{itemize}[leftmargin=*]
  \item \textbf{Verified student profile:} sign-up restricted to valid college email domains,
  with a structured profile listing tech stack, domain interests, experience level, and past
  project portfolio.
  \item \textbf{GitHub-linked verification:} students can link a GitHub profile/repository to
  their account, giving other students a verifiable signal of real project work alongside
  self-reported skills.
  \item \textbf{Team/project listings:} a student or existing team posts an open requirement
  (e.g., ``Need a Flutter developer for a fintech hackathon project, 48-hour build'') and
  specifies the roles/skills needed.
  \item \textbf{Matching engine:} recommends candidates by scoring skill complementarity,
  availability, and interest overlap, rather than plain keyword search.
  \item \textbf{Connection \& messaging:} in-app chat once a connection request is accepted, so
  students can coordinate without leaving the platform.
  \item \textbf{Team workspace:} once a team is formed, a shared space for task allocation and
  status tracking during the hackathon/project.
\end{itemize}

\subsection{Core workflow}
\begin{enumerate}[leftmargin=*]
  \item Student registers with a verified college email and builds a skill profile.
  \item Student either posts a requirement (``looking for teammates'') or browses/searches
  existing listings and profiles.
  \item The matching engine surfaces ranked candidate matches based on complementary skills.
  \item Both parties send/accept a connection request, unlocking in-app chat.
  \item If a team is formed, members get access to a shared workspace for the duration of the
  hackathon or project.
\end{enumerate}

\subsection{Operational constraints for an educational, multi-college setting}
\begin{itemize}[leftmargin=*]
  \item Keep the solution usable on low-to-mid range devices via responsive web design.
  \item Avoid dependence on advanced research algorithms; focus on workflow, verification, and
  matching correctness.
  \item Design for deployability using common web hosting and managed databases, so it can be
  piloted across multiple colleges without infrastructure overhead.
\end{itemize}

\section{Solution Approach \textit{...Engineering focus}}
This is an engineering course project, so we focus on scalable user, verification, and matching
workflows rather than novel algorithm research.

\subsection{Matching assistance \textit{...non-research}}
Teammate matching will be heuristic-based, not a research contribution:
\begin{itemize}[leftmargin=*]
  \item Normalize structured tags (tech stack, domain interest, experience level).
  \item Use a simple complementary-skill scoring measure (e.g., weighted overlap between
  ``skills offered'' and ``skills needed'') without claiming state-of-the-art recommendation
  performance. The score will combine skill complementarity with interest overlap, experience
  level, and availability using configurable weights.
  \item Show top-ranked candidate matches to the user for their own decision; never
  auto-assign teams.
\end{itemize}

\subsection{Web architecture \textit{...web-based solution}}
A typical 3-tier web architecture:
\begin{itemize}[leftmargin=*]
  \item \textbf{Frontend:} responsive UI for profile creation, browsing listings, and chat
  (React.js + Tailwind CSS).
  \item \textbf{Backend API:} authentication, profile/listing CRUD, matching logic, and
  notifications (Node.js + Express.js).
  \item \textbf{Database:} stores users, listings, connections, and messages (MongoDB Atlas).
\end{itemize}

\subsection{CI/CD and fast delivery enabler}
CI/CD will be used to:
\begin{itemize}[leftmargin=*]
  \item Automatically build and run tests on every change.
  \item Produce repeatable deployment artifacts to staging.
  \item Enable safe and frequent releases of incremental features (e.g., adding the matching
  engine after the core profile/listing flow is stable).
\end{itemize}

\section{Evaluation Criterion \textit{...measurable and attributable}}
We define success using measurable metrics tied to the core workflow.

\subsection{Primary evaluation metric}
\textbf{Time-to-first-match (TTM):} time between a student posting a team requirement and
receiving the first candidate match that satisfies a predefined relevance threshold (a minimum
skill/role compatibility score, not simply the first candidate returned).
\begin{itemize}[leftmargin=*]
  \item \textbf{Target:} median TTM $\leq$ 24 hours during the pilot, subject to sufficient
  active users and listings.
  \item \textbf{Attribution:} measured from database timestamps (listing/profile creation vs.
  first match surfaced by the matching engine).
\end{itemize}

\subsection{Secondary evaluation metrics}
\begin{itemize}[leftmargin=*]
  \item \textbf{Connection-to-team conversion rate:} percentage of accepted connections that
  result in a confirmed team.
  \item \textbf{Cross-college match rate:} percentage of successful matches formed between
  students of \emph{different} colleges (the platform's core value proposition).
  \item \textbf{Profile completeness / verification rate:} percentage of sign-ups that complete
  college email verification and a full skill profile.
  \item \textbf{GitHub linkage rate:} percentage of profiles with a linked GitHub repository,
  used as a proxy for verified technical activity.
  \item \textbf{User clarity:} student-reported clarity of the matching and connection flow
  using a simple 1--5 feedback question.
  \item \textbf{Reliability:} availability of staging/production for login and profile browsing
  during pilot hours (e.g., $\geq$ 99\% uptime).
\end{itemize}

\subsection{Pilot validation plan \textit{...high-level}}
\begin{itemize}[leftmargin=*]
  \item Recruit 15--25 pilot users across at least two colleges around an upcoming hackathon.
  \item Compare metrics before/after within the iteration window by logging baseline estimates
  via short interviews (``how did you find teammates last time?'').
\end{itemize}

\section{Timeline}
The project will be executed over a 12-week window using weekly iterations, consistent with
the time-to-value approach described earlier.

\begin{center}
\begin{tabular}{|p{4.3cm}|p{1.6cm}|p{6.3cm}|}
\hline
\textbf{Phase} & \textbf{Weeks} & \textbf{Deliverable} \\
\hline
Requirement analysis \& literature review & 1--2 & Finalized problem statement, SRS draft \\
\hline
System design (DB schema, architecture, wireframes) & 3--4 & ER diagram, architecture diagram, UI mockups \\
\hline
Backend development (auth, profiles, listings) & 5--7 & Working REST APIs, DB integration \\
\hline
Frontend development & 6--8 & Responsive web app connected to backend \\
\hline
Matching engine \& team workspace & 8--9 & Complementary-skill matching, shared team workspace \\
\hline
Testing (unit, integration, pilot UAT) & 10 & Test reports, bug fixes \\
\hline
Deployment \& documentation & 11 & Hosted demo, final report \\
\hline
Buffer, presentation \& submission & 12 & Final proposal defense / demo \\
\hline
\end{tabular}
\end{center}

\noindent\textbf{Note:} Weeks 1--7 constitute the \textbf{Initial MVP} (Profile $\rightarrow$
Listing $\rightarrow$ Search $\rightarrow$ Connection $\rightarrow$ Chat). Weeks 8--9 begin
\textbf{Iteration 2} (Matching Engine $\rightarrow$ Team Workspace $\rightarrow$ Analytics), as
detailed in the Project Scope and Deliverables section.

\section{Expected Outcomes}
\begin{itemize}[leftmargin=*]
  \item A working prototype demonstrating end-to-end team formation across at least two
  colleges.
  \item A validated complementary-skill matching engine with measurable time-to-first-match on
  pilot data.
  \item A measurable cross-college match rate demonstrating the platform's ability to facilitate
  inter-college collaboration.
  \item A reusable, documented codebase (GitHub repository: \texttt{campus-sync}) that can be
  extended beyond the course submission.
  \item A basis for a potential real-world pilot around an upcoming inter-college hackathon.
\end{itemize}

\section{Scalability \textit{...with foresight}}
The proposition should scale in both theoretical and practical senses.
\begin{enumerate}[leftmargin=*]
  \item \textbf{Scalable (theoretically):} The system should support growth in number of
  colleges, users, and listings by:
  \begin{itemize}
    \item decoupling components (frontend/backend),
    \item enabling pagination and indexing for common queries (search by skill, college, status),
    \item using stateless API design.
  \end{itemize}
  \item \textbf{Operationally deployable (practically):} The system should run on common web
  hosting:
  \begin{itemize}
    \item managed database,
    \item platform-hosted or containerized deployment,
    \item CI/CD pipeline for staging/production, so onboarding a new college is a
    configuration change, not a re-architecture.
  \end{itemize}
\end{enumerate}

\section{Engine availability heuristic}
A mature set of ``engines'' (tooling/services) is available to implement this as a web-based
project:
\begin{itemize}[leftmargin=*]
  \item Web framework for REST APIs and reactive pages.
  \item Managed document database.
  \item Authentication approach (token-based, with email-domain verification and optional OAuth).
  \item Real-time messaging library for in-app chat.
  \item Test framework for unit/integration tests.
  \item CI runner and staging deployment target.
\end{itemize}
We will use these mature components to focus on workflow design, verification, matching
correctness, and measurement rather than inventing new infrastructure.

\section{Project scope and deliverables \textit{...iteration-friendly}}
\subsection{Initial deliverable \textit{...first iteration}}
\begin{itemize}[leftmargin=*]
  \item College-email verified sign-up and structured skill profile
  \item Listing creation: post a ``looking for teammates'' requirement
  \item Basic browse/search of listings and profiles
  \item Connection request + in-app chat once accepted
  \item Basic logging and timestamps for TTM measurement
  \item CI: build + backend unit tests + linting
  \item CD to staging with a single command or pipeline job
\end{itemize}

\subsection{Subsequent deliverables}
\begin{itemize}[leftmargin=*]
  \item Complementary-skill matching engine with ranked recommendations
  \item Team workspace with task allocation and status tracking
  \item Admin dashboard for pilot metrics (TTM, conversion rate, cross-college rate)
  \item Improved search/filtering and indexed queries for scale readiness
\end{itemize}

\section{Risks and mitigations}
\begin{itemize}[leftmargin=*]
  \item \textbf{Cold start (too few users to match):} run the pilot around a scheduled hackathon
  so demand and supply of teammates arrive at the same time; seed listings manually if needed.
  \item \textbf{Identity / trust:} restrict sign-up to verified college email domains; keep
  match decisions user-initiated, never automatic.
  \item \textbf{Cross-college onboarding friction:} start with two colleges rather than an open
  rollout, to keep verification and support manageable.
  \item \textbf{Evaluation measurement ambiguity:} instrument timestamps and define clear
  listing/match/connection states before development begins.
  \item \textbf{Operational issues in pilot:} use staging early; ensure CI/CD produces
  repeatable deployments.
\end{itemize}

\section{Summary}
This project proposes a deployable web platform, CampusSync, to reduce the friction students
face when finding teammates for hackathons and projects outside their own college. It meets the
course criteria by emphasizing time-to-value through rapid iterations, implementing CI/CD to
support frequent safe releases, and using measurable evaluation criteria (especially
time-to-first-match and cross-college match rate). It remains focused on engineering workflow,
verification, and scalability readiness rather than research-driven novelty.

\section{References}
\begin{enumerate}[leftmargin=*]
  \item P. Resnick and H. R. Varian, ``Recommender systems,'' \textit{Communications of the
  ACM}, vol. 40, no. 3, pp. 56--58, 1997.
  \item L. N. Lee and A. B. Lee, ``Automated team formation in collaborative learning
  environments: A skill-based approach,'' \textit{Journal of Educational Technology \& Society},
  vol. 22, no. 1, pp. 1--12, 2019.
  \item F. Ricci, L. Rokach, and B. Shapira, \textit{Recommender Systems Handbook}, 2nd ed.
  Springer, 2015.
  \item C. Ren, Y. Wang, and X. Li, ``A survey on online collaboration platforms for
  distributed teams,'' in \textit{Proceedings of the ACM Conference on Computer-Supported
  Cooperative Work}, 2020.
  \item I. Sommerville, \textit{Software Engineering}, 10th ed. Pearson Education, 2016.
\end{enumerate}

\end{document}
